\section{What Is Logic?}

Logic is a method for systematically evaluating arguments of various
types. It can be used flexibly as a tool for identifying the strengths
and weaknesses of the various perspectives in a debate; as a way to sort through
the options in a mystery or puzzle; or as the bedrock of a hypothetico-deductive
model for science, which compares experimental results to theoretical
expectations in order to revise the theory or to design addititional
experiments.\footnote{See \cite{Smi20}[p. 17] for a slightly longer description
of the so-called \emph{hypothetico-deductive} model for science.}

Although this course is a study in \emph{formal} logic, the main concepts are
best introduced informally. Only by understanding these concepts can we begin
construct the formal definitions that we will use as the basis of our two
logics. And only by understanding these concepts can we even hope to see the
point of the extensive formal apparatus we will build over the next many weeks.

\subsection{Arguments}

An \emph{argument} is a sequence of sentences. It comprises some \emph{premises}
and a \emph{conclusion}.

\begin{argument}
\item[] My shoes are either in the den or in the bedroom.
\item[] My shoes are not in the den.
\item[\therefore] My shoes are in the bedroom.
\end{argument}

The conclusion in this \emph{inference} (i.e. ``one-step'' argument) is
introduced with a special symbol \therefore, an \emph{inference marker},
that can be pronounced ``so''.

In real-world examples, the conclusion can occur anywhere.

\begin{argument}
\item[] My shoes are either in the den or in the bedroom. They must be in
  the bedroom, because they are not in the den.
\end{argument}

But when we turn to formal logic, conclusions will be found at the end.

Informal arguments will often use special words to indicate which sentences
are meant to be premises and which is the conclusion. \\ 

\begin{table}[h]
  %\centering
  \begin{center}
    \begin{tabular}{|c|c|}
      \hline
      \textbf{conclusion indicators} & \textbf{premise indicators} \\
      \hline
      \makecell{ \emph{so, therefore, hence, thus,} \\
        \emph{accordingly, consequently, ... } } &
      \emph{since, because, given that, ... } \\
      \hline
    \end{tabular}
    \label{tab:indicator-words}
  \end{center}
  \captionsetup{justification=centering}
  \caption{Words commonly used to indicate premises \\
    and conclusions in ordinary English arguments.}
\end{table}

The premises and conclusions of these informal arguments are \emph{sentences},
but not just any type of sentence will do. We are not interested in commands or
questions, but only \emph{declarative} sentences. These are the sentences that
we can assign a \emph{truth value}; sentences that we can call \emph{true}
or \emph{false}.

In arguments like the one outlined above, the conclusion is a \emph{consequence}
of the premises; whenever the premises are true, the conclusion is also true.
When the consequence of an argument follows in this way from its premises, we
will say that the argument is \emph{valid}. Otherwise, it is \emph{invalid}.

Although we will keep the discussion informal for the time being, note that we
can talk about arguments by representing their premises and conclusions using
capital letters (with or without subscripts). Using this convention, we can
quickly define this notion of consequence (\emph{entailment}):

\begin{definition}[Informal Entailment]
  \label{def:informal-entailment}
  The sentences $A_1,\dots,A_n$ \emph{entail} the sentence $B$ if and
  only if $B$ is never true when all of $A_1,\dots,A_n$ are true.\footnote{
  From now on we will use ``iff'' to mean ``if and only if''.}
\end{definition}

We need to flesh out ``never'' a bit more in our formal definition, but this
gives us a sense of what we mean by an valid argument and its components.

\begin{definition}[Informal Validity]
  \label{def:informal-validity}
  An argument is \emph{valid} iff its premises entail its conclusion. Otherwise,
  the argument is \emph{invalid}.
\end{definition}

An invalid argument can be detected by noting a \emph{counter-example}: an
example of an argument with the same form where the premises are all true but
the conclusion is false.

%Arguments can be valid in a number of different ways. These depend on the
%limits we place on the types of counter-examples that we will accept as
%invalidating.

%\textbf{nomologically valid:} must respect the laws of nature;
%
%\textbf{conceptually valid:} must respect conceptual relations between words;

%Instead of focusing on these or other more contingent notions of validity,

We will be focused on a somewhat narrow notion of \emph{formal validity}.

We understand \emph{formal} as pertaining to the \emph{form} of arguments.
\footnote{Not, for example, as a synonym for \emph{fancy}.} Many conceptually
valid\footnote{respect conceptual relations between words} and nomologically
valid\footnote{respect the laws of nature} arguments cannot be \emph{formally}
validated. As the presentation goes on, we will more and more center our notion
of \emph{logical validity} on formal validity. 

The replacement of sentences with letters in arguments known to be valid is
our first step to \emph{formal logic}.

The argument on the left is valid while the one on the right is not.
\footnote{Unless we make clear the premise that $X\eif Z$}

\twoitems{
  \begin{argument}
  \item[] Either X or Y.
  \item[] Not Y.
  \item[\therefore] X.
  \end{argument}
} {
  \begin{argument}
  \item[] Either X or Y.
  \item[] Not Y.
  \item[\therefore] Z.
  \end{argument}
}

Note that the dependence on the \emph{forms} of arguments is fundamental here.
We can have both valid arguments with untrue conclusions and invalid arguments
whose premises and conclusions are all true. Without knowing actual content
of the sentences X, Y and Z, we can assess the arguments below as valid, based
entirely on the meaning of the other words being used: \emph{and}, \emph{or}
and \emph{not}.

\twoitems{
  \begin{argument}
  \item[] Both X and Y.
  \item[\therefore] X.
  \end{argument}
} {
  \begin{argument}
  \item[] Either X or Y.
  \item[] Not Y.
  \item[\therefore] X.
  \end{argument}
}

Moreover, these arguments are simple--indeed, minimal. In both cases, the
conclusion is justified by exactly the premise(s) provided, no more and no
less. We can use simple, rock-solid entailments of this sort as the basis of
a system of rules for building larger arguments (\emph{proofs}).

%\begin{quote}
%  One way of demonstrating that an inferential leap from some premises to a
%  given conclusion is valid is by breaking down the big leap into smaller
%  inference steps, each one of which is clearly valid.
%  \cite{Smi20}[p29, channeling Aristotle]
%\end{quote}

In the following complex argument (sort of a ``proto-proof'') we use several
\emph{inference steps} (marked with \therefore), each depending on a simple
rule and one or more of the sentences that came before it, some of which are
themselves conclusions.

To make the argument into a proof, we want to decorate it a little bit more,
most conspicuously by adding \emph{line numbers} and \emph{justifications} for
each inference step.

\begin{center}
  \begin{tabular}{c c c}
    (1) & Either X or Y.         & assumed                       \\
    (2) & If X, then Z.          & assumed                       \\
    (3) & Not Y.                 & assumed                       \\
    (4) & If Z and X, then W.    & assumed                       \\
    (5) & \therefore X.          & from 1 by ``or''              \\
    (6) & \therefore Z.          & from 2 and 5 by ``if...then'' \\
    (7) & \therefore Z and X.    & from 6 and 5 by ``and''       \\
    (8) & \therefore W.          & from 4 and 7 by ``if...then''
  \end{tabular}
\end{center}

\cite{Smi20} has a more extensive example, borrowed from Lewis Carroll, that
helps to illustrate how rules for the simple inference steps can be put to use
in such a fully-justified proof. Consider the following informal argument that
he labels (\textbf{D}):

\begin{argument}
\item[\textbf{D}] Anyone who understands human nature is clever; every true
  poet can stir the heart; Shakepeare wrote \emph{Hamlet}; no one who does not
  understand human nature can stir the heart; no one who is not a true poet
  wrote \emph{Hamlet}; so Shakespeare is clever.
  \cite{Smi20}[p31]
\end{argument}

Just two rules are needed to prove Carroll's argument, specifically.

\twoitems{
  \begin{argument}
  \item[ ] \textbf{Rule U:}
  \item[ ] Any/every B is C.
  \item[ ] a is B.
  \item[\therefore] a is C.
  \end{argument}
}{
  \begin{argument}
  \item[ ] \textbf{Rule V:}
  \item[ ] No one who is not A is B.
  \item[\therefore] Any B is A.
  \end{argument}
}

Like before, we take each inference step one at a time and list the lines that
we used to justify it. But now we can also add specific rule to our
justifications, allowing inference steps only when such a rule is available
(from \cite{Smi20}[p31]).

\begin{center}
  \begin{tabular}{c c c}
    (1) & Anyone who understands human nature is clever. & assumed \\
    (2) & Every true poet can stir the heart. & assumed \\
    (3) & Shakepeare wrote \emph{Hamlet}. & assumed \\
    (4) & No one who doesn't understand human nature can stir the heart. &
    assumed\\
    (5) & No one who is not a true poet wrote \emph{Hamlet}. & assumed \\
    (6) &\therefore Anyone who wrote \emph{Hamlet} is a true poet.&from 5 by V\\
    (7) & \therefore Shakespeare is a true poet. & from 6 and 3 by U \\
    (8) & \therefore Shakespeare can stir the heart. & from 2 and 7 by U \\
    (9) & \therefore Anyone who can stir the heart understands human nature. &
    from 4 by V \\
    (10) & \therefore Shakepeare understands human nature.&from 9 and 8 by U\\
    (11) & \therefore Shakespeare is clever. & from 10 and 1 by U \\
  \end{tabular}
\end{center}

\subsection{Enthymemes, Sound Arguments, Valid Arguments}

Returning to more concrete examples, we can see that many arguments in ordinary
English have assumptions that are silent or otherwise taken for granted. So the
following argument about my dog is correct.

\begin{argument}
\item[ ]          It is a perfect day for fetch.
\item[\therefore] Bandit is happy.
\end{argument}

But it is missing an important premise. Only once sentence (*) has been added
do we have an argument that is deductively valid.

\begin{argument}
\item[ ]          It is a perfect day for fetch.
\item[(*)]        If it is a perfect day for fetch, then Bandit is happy.
\item[\therefore] Bandit is happy.
\end{argument}

\begin{definition}[Enthymeme]
  \label{def:enthymeme}
  An argument with unsaid assumptions is called an \emph{enthymeme}.
\end{definition}

Logic is exceptional for detecting invalid arguments but cannot help at all with
assessing the truth of premises.

\begin{quote}
  The important thing to remember is that validity is not about the
  \emph{actual} truth or falsity of the sentences in the argument. It is about
  whether it is \emph{possible} for all the premises to be true and the
  conclusion to be not true at the same time. (\cite[][p15]{Mag25})
\end{quote}

This important point becomes a little more clear when reflecting on the contrast
between \emph{valid} arguments and \emph{sound} arguments. A valid argument
is just a matter of following the rules.

\begin{definition}[Valid Argument]
  An argument is \emph{valid} iff all of the inference steps from its premises
  to its conclusion are allowed by the relevant rules of logic.
\end{definition}

But a sound argument requires that we do not make any untrue assumptions.

\begin{definition}[Sound Argument]
  An argument is \emph{sound} iff (i) it is valid; \textbf{and} (ii) all of
  its premises are true.
\end{definition}

\subsection{Inductive Arguments, Deductive Arguments}

Our formal logics will have plenty to say about valid arguments, but will remain
mostly mute on the subject of sound arguments. Additionally, formal logic is
limited to the assessment of \emph{deductive} arguments. It does not help with
the assessment of the many \emph{inductive} arguments that populate both
empirical sciences and everyday experience.

\begin{argument}
\item[] The sun has risen every day so far.
\item[\therefore] The sun will rise tomorrow.
\end{argument}

While this is a perfectly reasonable inference to make, it is not a deductivley
valid argument. What do we mean by this?

\begin{definition}[Deductively Valid Argument]
  A \emph{deductively valid argument} is an argument where the conclusion
  follows directly from the premises via only a series of deductively valid
  steps: \emph{inference steps} such as those informally demonstrated above,
  where the conclusion follows from the premise(s) in \emph{every situtation}.
  \label{def:deductively-valid}
\end{definition}

We will not discuss \emph{inductive arguments} here, except to note that they
are not deductively valid. But we will need to eventually get a handle on what
we mean by ``situation''. First, we consider what kinds of inference steps are
appropriate to a deductive system of sentences.

Beyond valid inference steps, we can also define some other properties of
sentences and relations between them:

\Needspace{8\baselineskip}
\begin{definition}[Some logical properties of sentences]
  Sentences and collections of sentences can be assigned labels based on their
  truth in some or all situations:
  
  \begin{enumerate}
  \item Sentences that are always true are \emph{necessary truths}.
  \item Sentences that are never true are \emph{necessary falsehoods}.
  \item Sentences that are neither necessary truths nor necessary falshoods are
    \emph{contingent sentences}; they are sometimes true and sometimes false.
  \item Pairs of sentences that have the same truth value in every situation
    are \emph{necessarily equivalent}.
  \item Collections of sentences that are all true together in some situation
    are \emph{jointly possible}; those that are not are \emph{jointly
    impossible}.
  \end{enumerate}
  
\end{definition}

Negative results play an important role in the processes of logical deduction.
Most blatantly, if a sentence is not true, then its negation must be. The fact
that this rule works exceptionlessly is an important part of classical logic. In
a similar vein, we repeat the following rule from \cite[p. 13]{Smi20}.

\begin{definition}[The Invalidity Principle]
  \label{def:invalidity-principle}
  An inference step with actually true premises and an actually false conclusion
  must be invalid.  
\end{definition}

Thus when we find a counter-example to an argument, the argument is not valid.

\subsection{Proof by Absurdity}

One ancient approach to logical proofs is to show that assuming the opposite of
a conclusion will lead to a contradiction. We will see how the basic treatment
of negation leans heavily on this inference when we start to explore the syntax
and semantics of actual logics next week.

\begin{definition}[Reductio Ad Absurdum]
  \label{def:raa}
  If $\ephi_1\dots\ephi_n$ are taken as premises and the addition of \spsi
  would result in a contradiction, then $\ephi_1\dots\ephi_n$ entail not-\epsi.
\end{definition}

People often struggle taking reductio ad absurdum (RAA) seriously when they are
first introduced to it. One naive reading would take the following to be an
acceptable argument that Elvis is still alive.

\begin{argument}
  \item[ ] John lives in Amsterdam.
  \item[ ] It is not the case that John lives in Amsterdam.
  \item[\therefore] Elvis is still alive.
\end{argument}

No one likes this result, not even proponents of RAA. The most compelling
complaint is that Elvis' ongoing survival likely has nothing at all to do with
John's residence in the Netherlands. Still, like democracy, this is the worst
outcome except for all the rest. Most logicians are willing to accept this
result, despite its oddity, simply because the simplicity offered by the
overall system is worth preserving.

And it's worth noting that these examples are arguably just edge cases to begin
with. Indeed, by our existing definition of RAA the premises \sphi and
\enot\sphi are already inconsistent. This has a bunch of implications which
suggest that this premise set simply never comes up for serious consideration.

\begin{enumerate}
\item ``John lives in Amsterdam'' entails ``It is not the case that John
does not live in Amsterdam''.
\item ``John lives in Amsterdam'' and ``John does not live in Amsterdam''
  together entail ``Elvis is not alive.''
\end{enumerate}

Nevertheless, we will see RAA used all the time in service of formal proofs.

\subsection{Symbolization and Argument Schemas}

So far we showed how to use capital letters as stand-ins for sentences (i.e.
as \emph{schematic variables} in logical arguments and described this is as
fundamental to formal logics. Now we will discuss this process--often called
\emph{symbolization} or \emph{translation}--in more detail.

As we pull apart arguments attempting to ascertain their form, a basic
distinction that emerges is between \emph{content expressions} and \emph{topic
neutral terms}. So far, we have represented the former using capital letters
($A, B, C, X, Y, Z$) and the latter with their ordinary English counterparts
(\emph{and, not, either...or, if...then}. The important thing to note is that
formal logic is concerned primarily with the behavior of topic-neutral terms,
abstracting away the content expressions into mostly inert symbols.
\emph{Logical validity}\footnote{and \emph{logical entailment}} is just
deductive validity by virtue of topic-neutral terms alone. This is what
\emph{purely logical schemas} represent.

Consider again the two mini-arguments presented above, which we will give simple
names for reference.

\twoitems{
  \begin{argument}
  \item[] \textbf{And Rule:}
  \item[] Both X and Y.
  \item[\therefore] Z.
  \end{argument}
} {
  \begin{argument}
  \item[] \textbf{Or Rule:}
  \item[] Either X or Y.
  \item[] Not Y.
  \item[\therefore] X.
  \end{argument}
}

Thinking of \textbf{And Rule} and \textbf{Or Rule} as \emph{proto-schemas}, we
can glean a few informal observations:

\begin{observation}[Properties of Argument Proto-Schemas]
  \label{obs:argument-proto-schemas}
  Argument schemas have some regular properties:
  \begin{outline}[enumerate]
  \1 Schemas are intended to be a somewhat ``loose fit'' with the English
    expressions that they are intended to symbolize. The same logical form
    might apply to a whole range of different (but related) sentences of
    ordinary English.
    \1 Schemas depend on the radical abstraction of content expressions.
    \2 Letters are used to symbolize the content expressions of the object
    language.
    \2 The meaning of the expressions that these symbols represent is dropped or
    simplified dramatically.
    \1 The choice of content symbol is relatively unimportant, we just need to
    keep straight ``which ones are which.''
  \end{outline}
\end{observation}

The last point is worth lingering on. We can paraphrase it by saying that
schemas depend on the equivalence of substitution instances. In our rules, we
will use greek letters (\ephi, \epsi, \echi, ...) to represent arbitrary
expressions (mostly sentences):

\twoitems{
  \begin{argument}
  \item[ ] \textbf{And Rule:}
  \item[ ] \sphi and \epsi
  \item[\therefore] \epsi
  \end{argument}
} {
  \begin{argument}
  \item[ ] \textbf{Or Rule:}
  \item[ ] \sphi or \epsi
  \item[ ] not \ephi
  \item[\therefore] \epsi
  \end{argument}
}

By choosing a metalanguage variable in this role, we are stressing and important
property of our schemas, namely that they can be applied to all of the sentences
that share the form in question. Recall that we can describe a argument using
sentence letters to reveal its deeper formal structure. For example, we can
stipulate that each of the following inference steps is valid in virtue of its
form alone. \\

\threeitems{
  \begin{argument}
  \item[] A
  \item[] If A, then C
  \item[\therefore] C
  \end{argument}
}{
  \begin{argument}
  \item[] A or B
  \item[] not-A
  \item[\therefore] B
  \end{argument}
}{
  \begin{argument}
  \item[] not-(A and B)
  \item[] A
  \item[\therefore] not-B
  \end{argument}
}

The logic of sentences that we define is called \emph{Truth-Functional Logic}
($TFL$) or \emph{Propositional Logic} ($PL$), depending on the author. We will
use $TFL$ but take the two names to refer to one and the same logic.

I mentioned that logic is concerned with arguments built out of declarative
sentences. However, there are good reasons to use a different name for the
so-called sentences of formal logic. I repeat here just a few, as listed by
\cite{Smi20}[pp54-55]

\begin{observation}[TFL sentences are not like sentences of English]
  \begin{enumerate}
  \item Some grammatical English declarative sentences lack an interpretation
    and thus cannot be assigned a truth-value based on their sense.
  \item Some grammatical English sentences have more than one interpretation.
  \item Some TFL sentences can be used to express more than one grammatical
    English sentence.
  \item Some English sentences cannot be fully interpreted without consideration
    of the context.
  \item Surely there is a sense in which Aristotle and Frege are making many of
    \textbf{the same} arguments that we will consider: arguments that we would
    surely be mistaken to take as sentences of ancient Greek, late nineteenth
    century German or contemporary English
  \item TFL sentences cannot be distinguished except in their truth or falsity.
  \end{enumerate}
\end{observation}

Nevertheless, we will not adopt a specialized term (such as \emph{proposition})
for TFL sentences, noting only that we are using \emph{(TFL) sentence} as a
technical term, not as a synonym for the ordinary English concept
\emph{sentence}.

Recall that by limiting ourselves to sentence letters, we are marking off an
explicit boundary in our logic between those arguments that depend in some
way on the ``content'' of their sentences and those that depend on form alone.

When we look at Predicate Logic (Quantificational Logic) later this semester,
we will see how to incorporate the contributions of names and properties so
that we can encode \textbf{some} of the contributions that (sub-sentential)
content makes to the rules of inference. But, for now, we will focus on just
sentences and the forms of the arguments that they appear in.

%\begin{quote}
%  At best, TFL will help us to understand arguments that are valid due to
%  their form. \cite[][p30]{Mag25}
%\end{quote}

In the examples above, we have some sentences that contain \emph{sentence
letters} themselves. For example: ``If A, then B''. When we use capital
letters to stand for sentences (or subsentences) in this way, we are assuming
that the sentences represented are somewhat simple and can be directly assessed
as either \emph{true} or \emph{false}. We will call sentences of this sort
\emph{atomic sentences}. They cannot be broken down into parts and we cannot
``peek'' inside them to support our inferences.

When we use letters to represent sentences in this way, we will often provide a
\emph{symbolization key} to help decode the arguments.

\begin{argument}
\item[A:] It is a perfect day for fetch. 
\item[C:] Bandit is happy.
\end{argument}

Of course, these devices are \emph{ad hoc} since TFL is a formal logic and
the actual contents of its sentences is immaterial to the analysis.

\subsection{Logical Connectives}

We can build complex sentences from simpler ones (including atomic sentences)
with the help of a few \emph{connectives}. As the name suggests, the
connectives of TFL are \emph{truth-functional}. In other words, we will
interpret the truth of a complex sentence as a function of the truth of its
component subsentences based on the meanings of its connectives and the order
that they are introduced.

Consider the simple one-place connective $\lnot$. We can understand this
connective as a function from a truth-value to its opposite. So, for example,
when $A$ is true, $\lnot A$ is false. We can summarize the basic TFL
connectives as in the following table.

\begin{table}[h]
\centering
\begin{tabular}{|l|l|l|}
\hline
\textbf{Function} & \textbf{Symbol} & \textbf{Paraphrase} \\
\hline
Negation & $\lnot$ & 'It is not the case that \dots' \\
Conjunction & $\land$ & 'Both \dots and \dots' \\
Disjunction & $\lor$ & 'Either \dots or \dots' \\
Conditional & $\rightarrow$ & 'If \dots then \dots' \\
Biconditional & $\leftrightarrow$ & '\dots if and only if \dots' \\
\hline
\end{tabular}
\caption{Basic TFL Connectives}
\label{tab:basic-tfl-connectives}
\end{table}

The symbolization of sentences is subtle and not always a simple preservation
of the surface forms found in the English version. Consider each of the
connectives in turn.

\begin{argument}
\item[R:] A recession is inevitable.
\item[P:] A recession can be prevented.  
\end{argument}

If $P$ is part of our symbolization key, we might want to symbolize $R$ as
$\lnot P$ to capture the relationship between the sentences, even though
\emph{not} and \emph{it is not the case that} do not occur anywhere in
the sentence. Here is where the concept of \emph{paraphrase} comes into
play. We will want to lean pretty heavily on the flexibility of paraphrase,
especially when dealing with words that have antonyms or other 'lexical'
negations with \emph{un-}, \emph{non-}, \emph{de-}, and so on. 

We see a similar separation of English forms and TFL forms in our
treatment of conjunctions. If we want to symbolize \emph{Barbara is a medical
  doctor and a philosopher}, we need to depart from the English form and treat
it as a conjunction of sentences ($M \land P$) instead of as one sentence
(\emph{Barbara M and P}, perhaps). That is, we need to mention Barbara in
each of the sentences, not just the once. Our symbolization key would be:

\begin{argument}
\item[M:] Barbara is a medical doctor.
\item[P:] Barbara is a philosopher.  
\end{argument}

Note again that the semantics of TFL is bare bones. We will bundle under
conjunction several English connectives than have slightly different
implications in their full meaning: \emph{and}, \emph{but}, and
\emph{although} in the defintion, but many more in practice. In addition to
collapsing these English connectives to one in TFL, we will also ignore the
implications of sentence order. 

\begin{argument}
\item[A:] John left.
\item[B:] Mary got upset.
\end{argument}

Although \emph{A and B} and \emph{B and A} have different implications
about the order of events--and perhaps even about cause--in TFL $A \land B$
says the exact same thing as $B \land A$. But the interaction of the operators
is important. The following have a different meaning in TFL, even though each
could be the correct symbolization of the English sentence schema \emph{It
  is not the case that A and B}.

\twoitems{$\lnot(A \land B)$}{$(\lnot A \land B)$}

We need to take care in the order we apply the connectives.

With these preliminaries out of the way, we can now provide informal definitions
for each of the basic TFL connectives by way of paraphrase. We will use
lower-case Greek letters \ephi,\epsi,\echi,\dots to represent arbirtary TFL
sentences.

\begin{definition}[Truth-Functional Connectives]
  Truth-functional logic (TFL) comprises some atomic sentences and a few logical
  connectives to assemble them into sentences and arguments:
  \begin{enumerate}
  \item A sentence can be symbolized as \enot\sphi if it can be paraphrased as
    ''It is not the case that \ephi.''
  \item A sentence can be symbolized as \ephi\eand\spsi if it can be paraphrased
    as ``Both \sphi and \epsi'', as ``\ephi, but \epsi'' or as ``Although
    \ephi, \epsi''.
  \item A sentence can be symbolized as \ephi\eor\spsi if it can be paraphrased
    as ``Either \sphi or \epsi''.
  \item A sentence can be symbolized as \ephi\eif\spsi if it can be paraphrased
    as ``If \sphi, then \epsi'' or as ``\spsi only if \ephi''.
  \item A sentence can be symbolized as \ephi\eiff\spsi if it can be paraphrased
    as ``\sphi if and only if \epsi''.
  \end{enumerate}
\end{definition}
  
%\begin{definition}[Negation]
%\end{definition}
%\begin{definition}[Conjunction]
%\end{definition}
%\begin{definition}[Disjunction]
%\end{definition}
%\begin{definition}[Conditional]
%\end{definition}
%\begin{definition}[Biconditional]
%\end{definition}

  
